% part: intuitionistic-logic
% chapter: propositions-as-types
% section: normalization

\documentclass[../../../include/open-logic-section]{subfiles}

\begin{document}

\olfileid{int}{pty}{nor}

\olsection{సాధారణీకరణ}

ఈ విభాగంలో, తగిన తగ్గింపు క్రమంతో ఏ నిగమనాన్నైనా సాధారణ నిగమనంగా తగ్గించవచ్చని నిరూపిస్తాం; దీన్ని \emph{సాధారణీకరణ లక్షణం} అంటారు. ప్రతిపాదనలు-రకాల అనురూపతను వాడతాం: ప్రతి నిరూపణ పదం సాధారణ రూపానికి తగ్గించబడుతుందని చూపుతాం; దానివల్ల సహజ నిగమన \tetoken{నిగమనాల}{!!{derivation}s} సాధారణీకరణ అనుసరిస్తుంది.

ముందుగా పదాల సంక్లిష్టతను కొలిచే కొన్ని ప్రమేయాలను నిర్వచిద్దాం. \tetoken{ఒక సూత్రం}{!!a{formula}s} యొక్క \emph{పొడవు}~$\len{!A}$ నిర్వచనం:
\begin{align*}
  \len{p} &= 0\\
  \len{!A \land !B} &= \len{!A} + \len{!B} + 1\\
  \len{!A \lor !B} &= \len{!A} + \len{!B} + 1 \\
  \len{!A \lif !B} &= \len{!A} + \len{!B} + 1.
\end{align*}
రెడెక్స్~$M$ సంక్లిష్టతను దాని \emph{కట్ స్థాయి}~$\cutrank{M}$తో కొలుస్తాం:
\begin{align*}
  \cutrank{(\lambd[\typeof{x}{!A}][\typeof{N}{!B}])Q} &=
  \len{!A} + \len{!B} + 1 \\
  \cutrank{\ande{i}{\andi{\typeof{M}{!A}}{\typeof{N}{!B}}}} &=
  \len{!A} + \len{!B} + 1 \\
  \cutrank{\ore{\ori{i}{!A_{i}}{\typeof{M}{!A_i}}}
    {\typeof{x_1}{!A_1}}{\typeof{N_1}{!C}}
    {\typeof{x_2}{!A_2}}{\typeof{N_2}{!C}}} &=
  \len{!A_1} + \len{!A_2} + 1
\end{align*}
నిరూపణ పదం సంక్లిష్టత దానిలోని అత్యంత సంక్లిష్ట రెడెక్స్ స్థాయి; పదం సాధారణమైతే $0$:
\[
  \maxrank{M} = \max (\{0\} \cup \{\cutrank{N} | N \text{ అనేది } M \text{ యొక్క రెడెక్స్ ఉపపదం}\})
\]
\sourcecorrection{OLTEPTPTYNOR-001}{వియోగ రెడెక్స్ రకాలు $!A_1$, $!A_2$గా ఉన్నా, మూలం స్థాయి సూత్రంలో $!A$, $!B$లను వాడింది; ప్రకటించిన రకాలతో సమన్వయం చేశాం.}
\sourcecorrection{OLTEPTPTYNOR-002}{సాధారణ పదానికి స్థాయి $0$ అని గద్యం చెబుతున్నా, మూల సూత్రం ఖాళీ సమితి గరిష్ఠాన్ని తీసుకుంటుంది. గరిష్ఠం తీసే సమితిలో $0$ను చేర్చాం.}

\begin{lem}\ollabel{lem:subst}
  $\Subst{M}{\typeof{N}{!A}}{\typeof{x}{!A}}$ రెడెక్స్ అయి, $M
  \not\ident x$ అయితే, కింది సందర్భాలలో ఒకటి వర్తిస్తుంది:
  \begin{enumerate}
  \item $M$ స్వయంగా రెడెక్స్; లేదా
  \item $M$ రూపం $\ande{i}{x}$, $N$ రూపం $\andi{P_1}{P_2}$;
  \item $M$ రూపం $\ore{x}{x_1}{P_1}{x_2}{P_2}$, $N$ రూపం $\ori{i}{}{Q}$;
  \item $M$ రూపం $x Q$, $N$ రూపం $\lambd[x][P]$.
  \end{enumerate}
  మొదటి సందర్భంలో $\cutrank{\Subst{M}{N}{x}} = \cutrank{M}$; ఇతర సందర్భాల్లో $\cutrank{\Subst{M}{N}{x}} = \len{!A}$.
\end{lem}
\begin{proof}
  $M$పై ఆగమనంతో నిరూపిద్దాం.
  \begin{enumerate}
  \item $M$ ఒకే చరం~$y$ అయి, $y \not\ident x$ అయితే, $\Subst{y}{N}{x}$ అనేది~$y$; కాబట్టి రెడెక్స్ కాదు.
  \item $M$ రూపం $\andi{N_1}{N_2}$, $\lambd[x][N]$, లేదా $\ori{i}{!A}{N}$ అయితే, $\Subst{M}{\typeof{N}{!A}}{\typeof{x}{!A}}$ కూడా అదే రూపంలో ఉంటుంది; కాబట్టి రెడెక్స్ కాదు.
  \item $M$ రూపం $\ande{i}{P}$ అయితే, రెండు సందర్భాలను పరిగణిద్దాం.
    \begin{enumerate}
      \item $P$ రూపం $\andi{P_1}{P_2}$ అయితే, $M \ident
        \ande{i}{\andi{P_1}{P_2}}$ రెడెక్స్; స్పష్టంగా
        \[
        \Subst{M}{N}{x} \ident
        \ande{i}{\andi{\Subst{P_1}{N}{x}}{\Subst{P_2}{N}{x}}}
        \]
        కూడా రెడెక్స్. కట్ స్థాయిలు సమానం.
      \item $P$ ఒకే చరం అయితే, ప్రతిస్థాపన రెడెక్స్ కావాలంటే అది~$x$ కావాలి; $N$ రూపం $\andi{P_1}{P_2}$ కావాలి. ఇప్పుడు
        \[
        \Subst{M}{N}{x} \ident \Subst{\ande{i}{x}}{\andi{P_1}{P_2}}{x},
        \]
        అనేది $\ande{i}{\andi{P_1}{P_2}}$. దాని కట్ స్థాయి ప్రతిస్థాపించే చరం రకం యొక్క పొడవు~$\len{!A}$.
    \end{enumerate}
  \end{enumerate}
  $\ore{N}{x_1}{N_1}{x_2}{N_2}$, $P Q$ సందర్భాలు ఇలాంటివే.
\end{proof}
\sourcecorrection{OLTEPTPTYNOR-003}{లెమ్మాలో వియోగ తొలగింపు పదం ప్రధాన స్థానంలో మూలం $i$ను రాసింది; ప్రతిస్థాపించే చరం $x$ అక్కడ ఉండాలి.}
\sourcecorrection{OLTEPTPTYNOR-004}{లెమ్మా చివరి స్థాయి సమీకరణంలో అవసరం లేని మూసే కుండలీకరణాన్ని తొలగించాం.}
\sourcecorrection{OLTEPTPTYNOR-005}{మూలం చరం $x$కు $\cutrank{x}$ ఉందని చెప్పింది; కట్ స్థాయి రెడెక్స్‌కు మాత్రమే నిర్వచించబడింది. ఉద్దేశించినది $x$ రకం పొడవు $\len{!A}$.}

\begin{lem}
  $M$ అనేది $M'$కు సంకోచితం అయి, $M$లోని ప్రతి సరైన రెడెక్స్ ఉపపదం~$N$కు $\cutrank{M} > \cutrank{N}$ అయితే, $\cutrank{M} >
  \maxrank{M'}$.
\end{lem}

\begin{proof}
  సందర్భాలవారీగా నిరూపిద్దాం.
  \begin{enumerate}
  \item $M$ రూపం $\ande{i}{\andi{M_1}{M_2}}$ అయితే, $M'$ అనేది $M_i$; $M_i$లోని ప్రతి ఉపపదమూ $M$కు సరైన ఉపపదమే కాబట్టి ఫలితం వస్తుంది.
  \item $M$ రూపం $(\lambd[\typeof{x}{!A}][N])\typeof{Q}{!A}$ అయితే, $M'$ అనేది $\Subst{N}{\typeof{Q}{!A}}{\typeof{x}{!A}}$. $M'$లోని ఒక రెడెక్స్‌ను పరిగణించండి. అది $N$ లేదా $Q$లోని సమాన కట్ స్థాయి గల రెడెక్స్‌కు అనురూపం కావచ్చు; ఆ స్థాయి ఉపపత్తి ప్రకారం $\cutrank{M}$ కంటే తక్కువ. లేదా కట్ స్థాయి $\len{!A}$ అవుతుంది; నిర్వచనం ప్రకారం ఇది $\cutrank{(\lambd[x^{!A}][N])Q}$ కంటే తక్కువ.
  \item $M$ రూపం
    \[
    \ore{\ori{i}{}{\typeof{N}{!A_i}}}
        {\typeof{x_1}{!A_1}}{\typeof{N_1}{!C}}
        {\typeof{x_2}{!A_2}}{\typeof{N_2}{!C}},
    \]
    అయితే, $M' \ident \Subst{N_i}{N}{\typeof{x_i}{!A_i}}$. ~$M'$లోని ఒక రెడెక్స్‌ను పరిగణించండి. అది $N_i$ లేదా $N$లోని సమాన కట్ స్థాయి గల రెడెక్స్‌కు అనురూపం కావచ్చు; ఆ స్థాయి ఉపపత్తి ప్రకారం $\cutrank{M}$ కంటే తక్కువ. లేదా కట్ స్థాయి $\len{!A_i}$ అవుతుంది; నిర్వచనం ప్రకారం ఇది $\cutrank{\ore{\ori{i}{}{\typeof{N}{!A_i}}}
      {\typeof{x_1}{!A_1}}{\typeof{N_1}{!C}}
      {\typeof{x_2}{!A_2}}{\typeof{N_2}{!C}}}$ కంటే తక్కువ.
  \end{enumerate}
\end{proof}
\sourcecorrection{OLTEPTPTYNOR-006}{మూలం ప్రతిస్థాపన జరిగే పదంలోని పాత రెడెక్స్‌లు, కొత్తగా ఏర్పడే రెడెక్స్‌లను మాత్రమే పేర్కొంది; చేర్చే పదం నుంచి కాపీ అయ్యే రెడెక్స్‌లను కూడా పరిగణించాలి. అవీ అసలు రెడెక్స్‌కు సరైన ఉపపదాల్లోనే ఉన్నాయి కనుక ప్రకటించిన తక్కువ-స్థాయి ఉపపత్తి వాటికీ వర్తిస్తుంది.}

\begin{thm}
  అన్ని నిరూపణ పదాలూ సాధారణ రూపానికి తగ్గించబడతాయి; అన్ని \tetoken{నిగమనాలూ}{!!{derivation}s} సాధారణ \tetoken{నిగమనాలుగా}{!!{derivation}s} తగ్గించబడతాయి.
\end{thm}

\begin{proof}
  రెండవది మొదటిదాని నుంచి వస్తుంది. నిరూపణ పదం~$M$కు, $m=\maxrank{M}$పై సంపూర్ణ ఆగమనంతో మొదటిదాన్ని నిరూపిద్దాం.
  \begin{enumerate}

  \item $m=0$ అయితే, $M$ ఇప్పటికే సాధారణ రూపంలో ఉంది.
  \item
    లేకపోతే, $M$లో కట్ స్థాయి~$m$కు సమానంగా ఉన్న రెడెక్స్‌ల సంఖ్య~$n$పై ఆగమనం చేద్దాం.
    \begin{enumerate}
    \item $n=1$ అయితే, రెడెక్స్ అయిన ప్రతి సరైన ఉపపదం~$P$కు $m = \cutrank{N} >
      \cutrank{P}$ అయ్యే రెడెక్స్~$N$ను ఎంచుకోండి. ప్రతి పదానికీ పరిమితంగా ఉపపదాలు మాత్రమే ఉన్నాయి కాబట్టి అటువంటి రెడెక్స్ తప్పనిసరిగా ఉంటుంది.

      $N$ సంకోచన ఫలితాన్ని $N'$గా సూచిద్దాం. లెమ్మా ప్రకారం $\maxrank{N'} < \maxrank{N}$; కాబట్టి $\cutrank=m$ గల రెడెక్స్‌ల సంఖ్య~$n$ తగ్గుతుంది. అందువల్ల~$m$ తగ్గుతుంది ($1$ లేదా అంతకంటే ఎక్కువ); అప్పుడు~$m$పై ఆగమన ఉపపత్తిని వర్తింపజేయవచ్చు.
    \item ఆగమన దశలో $n>1$ అనుకుందాం. ప్రక్రియ ఇలాంటిదే; అయితే~$n$ తగ్గినా ధన సంఖ్యగా మిగులుతుంది,~$m$ మారదు. ~$n$పై ఆగమన ఉపపత్తిని వర్తింపజేస్తాం.
    \end{enumerate}
\end{enumerate}
\end{proof}
\sourcecorrection{OLTEPTPTYNOR-007}{ఈ మూల వాదనలో స్థానిక రెడెక్స్ స్థాయి తగ్గుదల నుంచి మొత్తం పదంలోని గరిష్ఠ స్థాయి రెడెక్స్‌ల సంఖ్య తగ్గిందని నేరుగా నిర్ధారించారు. చుట్టూ కొత్త రెడెక్స్ ఏర్పడే సందర్భాన్ని ఇది సమర్థించదు. ఉదాహరణకు తక్కువ స్థాయి వియోగ రెడెక్స్ ఫలితం అధిక సంక్లిష్ట లబ్ధ రకం అయితే, దాని పైనున్న ప్రక్షేపణ కొత్త అధిక-స్థాయి రెడెక్స్ అవుతుంది. కాబట్టి ఈ కొలమాన వాదన పూర్తి సాధారణీకరణ నిరూపణ కాదు; తరువాత ప్రకటించే బలమైన సాధారణీకరణకూ వేరే వాదన అవసరం. మూల వాదనను పరిశీలన కోసం నిలిపాం.}

రకాలతో కూడిన $\lambd$-పదాలు సాధారణీకరించబడతాయంటే, ప్రతి పదానికి \emph{సాధారణ రూపం ఉంది}. $\lambd$-పదాలను $\beta$-తగ్గింపుతో సాధారణ రూపానికి తేవడాన్ని గణన అమలుగా, సాధారణ రూపాన్ని గణన విలువగా భావిస్తే, రక $\lambd$-కలనశాస్త్రంలోని ప్రతి గణన చివరకు ఆగి విలువను ఇస్తుందని అర్థం. రక సంరక్షణ వల్ల ఆ విలువకు సరైన రకం ఉంటుంది. ఉదాహరణకు $N$ రకం $!A \lif !B$ అయితే (అంటే అది~$!A$ రకం ఆర్గ్యుమెంట్లను తీసుకుని~$!B$ రకం విలువలను ఇచ్చే ప్రమేయం), దాన్ని~$!A$ రకం పదం~$M$కు (ఇన్‌పుట్) ప్రయోగిస్తే, $(NM)$ పదం~$N'$కు (అవుట్‌పుట్) తగ్గుతుంది; ఆ ఫలితం రకం~$!B$ అని హామీ ఉంటుంది.

నిజానికి, రక $\lambd$-కలనశాస్త్రంలోని పదాలు ఉపపదాలకు $\beta$-తగ్గింపులు వర్తింపజేసే ఒకే క్రమంలో సాధారణ రూపాలకు తగ్గడమే కాదు; $\beta$-తగ్గింపులను వర్తింపజేసే \emph{ఏ} క్రమమైనా చివరకు సాధారణ రూపంలో ముగుస్తుంది. ఈ లక్షణాన్ని \emph{బలమైన సాధారణీకరణ} అంటారు. అంటే గణన ఆగి విలువను ఇస్తుందని నిర్ధారించే ఒక విధానం ఉండటమే కాదు; గణనను అమలుచేసే \emph{ఏ} విధానమైనా చివరకు విలువను ఇస్తుంది.

గణనను వేర్వేరు విధాలుగా అమలుచేసినా \emph{అదే} విలువ వస్తుందని ఎలా తెలుసుకోగలం? ఇప్పటివరకు రక సంరక్షణ వల్ల అన్ని ఫలిత విలువలకూ ఒకే రకం ఉంటుందని మాత్రమే తెలుసు. రక $\lambd$-కలనశాస్త్రానికి మరో ముఖ్య లక్షణం ఉంది: అది \emph{సంగమ} లేదా \emph{చర్చ్--రోసర్} లక్షణాన్ని తీరుస్తుంది. $N \red N_1$, $N \red
N_2$ అయితే, $N_1 \red N'$, $N_2 \red
N'$ అయ్యే పదం~$N'$ ఉంటుంది. $N \red N'$ అంటే $N'$ అనేది సున్నా లేదా ఎక్కువ $\redone$-తగ్గింపు దశల ఫలితమని గుర్తుంచుకోండి. $N_1$ సాధారణ రూపంలో ఉంటే, $N_1 \red N'$ అయ్యే ఏకైక పదం~$N'$, $N_1$ స్వయమే (సున్నా $\redone$-దశలతో). కాబట్టి $N_1$, $N_2$ రెండూ సాధారణ రూపాల్లో ఉంటే $N_1 = N' = N_2$. అంటే $\redone$ బలంగా సాధారణీకరించబడుతుంది కాబట్టి ప్రతి తగ్గింపు క్రమం చివరకు సాధారణ రూపాన్ని ఇస్తుంది. $\red$కు చర్చ్--రోసర్ లక్షణం ఉన్నందున ఏ రెండు సాధారణ రూపాలైనా సమానం. అందువల్ల రక $\lambd$-కలనశాస్త్రంలోని గణన ఎప్పుడూ ముగుస్తుంది; గణన ఏ క్రమంలో సాగినా విలువ (సాధారణ రూపం) ఒకటే.
\sourcecorrection{OLTEPTPTYNOR-010}{మూలం బలమైన సాధారణీకరణను స్వప్రావర్తక-సంక్రమణ బంధమైన $\red$కు ఆపాదించింది. సున్నా దశలను అనుమతించే ఆ సంబంధానికి అనంత స్వీయ దశలు ఉంటాయి; బలమైన సాధారణీకరణ వర్తించేది ఒక్కదశ సంబంధం $\redone$కు. సంబంధిత వాక్యాలు, న్యూమన్ లెమ్మాలో ఆ గుర్తును సరిచేశాం.}

సంబంధానికి సంగమ లక్షణం ఉందని నిరూపించడం సాధారణంగా కష్టం. అయితే కింది బలహీన షరతును సులభంగా స్థాపించవచ్చు:

\begin{prop}[బలహీన చర్చ్--రోసర్ లక్షణం]
$N \redone N_1$, $N \redone N_2$ అయితే, $N_1 \red N'$, $N_2 \red N'$ అయ్యే పదం~$N'$ ఉంది.
\end{prop}

\begin{proof}
  $N$లోని రెడెక్స్~$M_1$ లేదా~$M_2$ను దాని సంకోచన ఫలితంతో మార్చితే వరుసగా $N_1$, $N_2$ వస్తాయి. రెడెక్స్‌లు $M_1$, $M_2$ అనేవి $N$ ఉపపదాలు; కాబట్టి అవి పాక్షికంగా ఒకదానితో ఒకటి కలవలేవు. అంటే ఒకటి మరొకదాని ఉపపదమై ఉండాలి, లేదా అవి $N$లో వేర్వేరు, పరస్పరం కలవని స్థానాల్లో ఉండాలి. రెండవ సందర్భంలో $N$ రూపం $N(M_1,M_2)$. $M_1$ స్థానంలో దాని సంకోచన ఫలితం~$M_1'$ను ఉంచితే $N_1$ రూపం $N(M_1',M_2)$; $M_2$ స్థానంలో దాని సంకోచన ఫలితం~$M_2'$ను ఉంచితే $N_2$ రూపం $N(M_1,M_2')$. రెండూ $N(M_1',M_2')$కు తగ్గుతాయి. ఇతర సందర్భంలో $M_2$ అనేది~$M_1$ ఉపపదమని అనుకుందాం (మరో సందర్భం సమమితమైనది); $M_1$ ఏ రకపు రెడెక్స్ అన్నదాని ప్రకారం సందర్భాలను వేరు చేద్దాం. ఉదాహరణకు $M_1 \ident
  \proj{1}{\pair{O_1}{O_2}}$ అయితే అది $O_1$కు తగ్గుతుంది. $M_2$ అనేది $O_1$ ఉపపదమైతే, $M_2$ను పరివర్తితం చేస్తే $O_1'$ వస్తుంది. కాబట్టి ముందు $M_1$ను, తరువాత $M_2$ను పరివర్తితం చేస్తే $O_1'$ వస్తుంది. కానీ ముందు $M_2$ను పరివర్తితం చేస్తే $\proj{1}{\pair{O_1'}{O_2}}$ వస్తుంది; ఇదీ~$O_1'$కు తగ్గుతుంది.
\end{proof}
\sourcecorrection{OLTEPTPTYNOR-008}{మూలంలోని చివరి ప్రక్షేపణ ఫలితం $O_2$గా ఉంది; మొదటి భాగ ప్రక్షేపణ ఫలితం $O_1'$ కావాలి.}
\sourcecorrection{OLTEPTPTYNOR-009}{ఈ మూల నిరూపణ పరస్పరం వేర్వేరు రెడెక్స్ స్థానాలను, ఒక ప్రక్షేపణ అంతర్భాగ సందర్భాన్ని మాత్రమే చూపిస్తుంది. ఇతర అంతర్భాగ స్థానాలు, ప్రయోగ-వియోగ రెడెక్స్‌లు, ప్రతిస్థాపనతో కలిగే ప్రతులు/తొలగింపుల సంగమ వాదనలు ఇంకా ఇవ్వలేదు. మూల ఉదాహరణను నిలిపాం; ఇది పూర్తి బలహీన చర్చ్--రోసర్ నిరూపణ అని ధృవీకరించలేదు.}

\begin{lem}[న్యూమన్ లెమ్మా]
$\redone$ బలంగా సాధారణీకరించబడి, బలహీన చర్చ్--రోసర్ లక్షణాన్ని తీరుస్తే, దాని బహుళదశ తగ్గింపు సంబంధం సంగమ లక్షణాన్ని తీరుస్తుంది.
\end{lem}

\begin{proof}
  $\redone$ బలంగా సాధారణీకరించబడుతుంది కాబట్టి, ప్రతి తగ్గింపు వరుస సాధారణ రూపంలో ముగుస్తుంది. $\red$ సంగమ లక్షణాన్ని తీరుస్తేనే సాధారణ రూపాలు అనన్యంగా ఉంటాయి. కాబట్టి పదం~$N$కు వేర్వేరు సాధారణ రూపాలు~$N'$, $N''$ ఉన్నాయని అనుకుందాం; అవి కింది తగ్గింపు వరుసల ఫలితాలు:
  \begin{align*}
  N & \redone N_1' \redone N_2' \redone \dots \redone N_k' = N' \text{ మరియు}\\
  N & \redone N_1'' \redone N_2'' \redone \dots \redone N_l'' = N''
  \end{align*}
  (తగ్గింపు వరుస పొడవు $>0$ కావాలి; లేకపోతే~$N$ ఇప్పటికే సాధారణ రూపంలో ఉంటుంది.)
  \sourcecorrection{OLTEPTPTYNOR-011}{మూల వరుసల్లో $N=N_1'=N_1''$ అని మొదలుపెట్టి, తరువాత $N_1'$, $N_1''$లను మొదటి వేర్వేరు ఫలితాలుగా వాడింది. ఆ ఉద్దేశానికి అనుగుణంగా ప్రారంభం నుంచి మొదటి ఫలితానికి ఒక్కదశ తగ్గింపును రాశాం.}

  $N_1' = N_1''$ అయితే ఆ పదానికి రెండు వేర్వేరు సాధారణ రూపాలు ($N'$, $N''$) ఉన్నాయి. $N_1' \neq N_1''$ అయితే బలహీన చర్చ్--రోసర్ వల్ల $N_1' \red N'''$, $N_1'' \red N'''$ అయ్యే పదం~$N'''$ ఉంది. బలమైన సాధారణీకరణతో ఆ ఉమ్మడి పదాన్ని సాధారణ రూపం~$N^\ast$కు తగ్గించండి: $N''' \red N^\ast$. $N' \neq N''$ కాబట్టి, $N^\ast \neq N'$ లేదా $N^\ast \neq N''$. అందువల్ల $N_1'$ లేదా $N_1''$కు రెండు వేర్వేరు సాధారణ రూపాలు ఉన్నాయి. అంటే~$N$కు రెండు వేర్వేరు సాధారణ రూపాలు ఉంటే, $N \redone N_1$ అయ్యే పదం~$N_1$ ఉండి, $N_1$కూ రెండు వేర్వేరు సాధారణ రూపాలు ఉంటాయని చూపాం. దీన్ని కొనసాగిస్తే $N \redone N_1
  \redone N_2 \redone \dots$ అనే అనంత వరుస వస్తుంది; కానీ $\redone$ బలంగా సాధారణీకరించబడుతుంది కాబట్టి ఇది అసాధ్యం.
  \sourcecorrection{OLTEPTPTYNOR-012}{మూలం ఉమ్మడి ఫలితం $N'''$ను సాధారణ రూపమని నిర్ధారించకుండానే దాన్ని $N'$, $N''$లతో పోల్చి రెండు వేర్వేరు సాధారణ రూపాలున్నాయని తేల్చింది. ముందుగా ఉమ్మడి ఫలితాన్ని సాధారణ రూపానికి తగ్గించే దశను చేర్చాం.}
\end{proof}

\end{document}
